\chapter*{前言}


1980年，彼得·德鲁克写下《动荡时代的管理》。他在那个冷战未歇、滞胀持续的年代提出了一个核心警告：动荡时代最大的危险不是动荡本身，而是沿用过去的逻辑做事。

四十多年后，这句话的穿透力有增无减。只不过，驱动这场动荡的不再是地缘政治或石油危机，而是一场技术革命——人工智能。

本书无意重写德鲁克的著作。德鲁克的理论框架属于他的时代，他的分析方法却属于所有时代。本书所做的，是借他的目光审视我们这个时代，追问一个核心问题：

\textbf{当AI重新定义组织的运作逻辑时，我们该如何重建管理的根基？}

老王我年少的时候，一度轻狂的认为管理这事挺简单的——定战略、设目标、开例会、追进度。后来慢慢才发现，管人比管事难一百倍。管事有流程，管人没剧本。

这些年在互联网圈混，见过太多起起落落。有人追到了风口，飞起来了；有人踩中了浪尖，结果浪退了才发现自己在裸泳。我自己呢？跟风过几个概念，几乎没有成功的。

但最让我感慨的不是成功或失败本身，而是——每一波浪潮来的时候，大家都在讨论"这次有什么不一样"，可潮水退了才发现，\textbf{真正决定胜负的，从来不是技术本身，而是组织怎么去使用技术}。


2016年深度学习火的时候，所有人都在谈算法、谈算力、谈数据。几年后回头看，活下来的公司不是因为算法更好，而是组织能力更强。2022年大模型来了，又一轮"AI颠覆一切"的浪潮。这次我学乖了，先不问"AI能做什么"，先问"我们这个组织，能不能接得住"。

有个现象挺有意思的：同样的AI项目，不同组织来做，时间线能差好几倍。有人三个月上线，有人三年还在试点。差距在哪儿？不是技术，是组织——是决策权怎么分、流程怎么改、人怎么用。翻译成人话就是：你还在那儿磨蹭，人家已经跑出去三条街了。

德鲁克1980年写《动荡时代的管理》时，互联网还没商业化。他不可能预见到AI。但他的方法论——从本质出发，追问组织存在的根本理由——让他说的每句话，放到今天都像是直接在说AI时代的事。

所以这本书，不是讲AI的书。这本书讲的是：\textbf{AI来了，做管理的人，脑子里那套框架需要怎么换？}

这不是一个技术问题。这是一个关于人的问题，关于组织的问题，关于"在变化中怎么不变"的问题。

本书分为七个部分：

\begin{itemize}
\item \textbf{时代背景}——为什么AI不是又一个技术浪潮，而是文明级的转折
\item \textbf{思维重构}——从预测到塑造，从控制到涌现
\item \textbf{流程重构}——控制逻辑失效后，流程如何重新设计
\item \textbf{组织重构}——从机械系统到生物系统
\item \textbf{技能重构}——可编码技能的贬值与不可编码技能的升值
\item \textbf{阻力维度}——变革真正的障碍来自哪里
\item \textbf{面向未来}——德鲁克如果今天重写，他会改变什么

\end{itemize}
感谢你翻开这本书。如果它能让你在某个问题上看得稍微通透一些，那便是它存在的意义。

如果不能，那就当老王我又开始不懂装懂絮絮叨叨地瞎聊易经，就当脱口秀看吧。

\textit{2026年春}
